% ===========================================================================
% Abstract.
% Aufbau und Position (nach "Verwendung von KI-Werkzeugen", vor dem
% Inhaltsverzeichnis) nach der Vorlage in
% ../master_thesis_obsidian/literature/jonas_master_thesis.pdf (S. v).
% Alle Kennzahlen entsprechen Kapitel 6 (Kernstichprobe, GPT-5.5).
% ===========================================================================
\clearpage
\thispagestyle{plain}
\setheader{Abstract}

\chapter*{Abstract}

\noindent
Das nationale Metadatenportal GovData macht offene Verwaltungsdaten von Bund, Ländern
und Kommunen über Metadaten nach dem Standard DCAT-AP.de auffindbar. Deren Qualität
wird bislang vor allem formal über eine SHACL-Validierung geprüft. Eine solche Prüfung
sagt aus, ob ein Metadatensatz standardkonform ist, nicht aber, ob er für Nutzende
verständlich, erreichbar und nachnutzbar ist.

Die vorliegende Arbeit entwickelt ein Bewertungsmodell, das formale Prüfungen um
technische sowie kontext- und semantikbezogene Indikatoren ergänzt, setzt es
prototypisch um und evaluiert es. Aus einer Literaturanalyse und drei
Experteninterviews werden fünfzehn Kandidatendimensionen abgeleitet und zu den vier
Dimensionen Auffindbarkeit, Zugänglichkeit, Nachnutzbarkeit und Aussagekraft
verdichtet, die über 27 Einzelindikatoren operationalisiert sind. Sechs davon bewerten
die Aussagekraft der Freitextfelder mithilfe eines Sprachmodells. Ein menschenlesbares
Reporting überführt jedes negative Prüfergebnis in einen Befund mit Fundstelle und
Handlungsanweisung.

Evaluiert wird auf einer stratifizierten Zufallsstichprobe von 50 GovData-Datensätzen
sowie auf 20 konstruierten Extremfällen, jeweils gegen ein manuelles Referenzurteil
und gegen eine Reimplementierung des Metadata Quality Assessment (MQA) von
data.europa.eu. Auf der Stichprobe erreicht der Prototyp eine Rangkorrelation von
$\rho = 0{,}765$ mit dem Referenzurteil und trifft in 92\,\% der Fälle dieselbe
Gesamteinordnung. Die MQA nimmt auf tiefer geprüften Indikatoren in 15,8\,\% der Fälle
eine nicht gegebene Qualität an, und gegen das Referenzurteil liegt der Prototyp auf
der Stichprobe bei fünf von sechs Zielgrößen vorn. Schwächste Dimension ist die
sprachmodellgestützte Aussagekraft. Einschränkend stammt das Referenzurteil von einer
einzelnen Person, und die Gewichte des Modells sind nicht empirisch validiert.
Quellcode und Auswertungen sind öffentlich verfügbar.\footnote{\url{https://github.com/LarsBuergerr/master_thesis_prototype}}

\pagebreak